\documentclass[11pt,a4paper,fontset=fandol]{ctexart}

\usepackage[a4paper,top=2.35cm,bottom=2.55cm,left=2.3cm,right=2.3cm,headheight=18pt,headsep=0.55cm,footskip=26pt]{geometry}
\usepackage{amsmath,amssymb,amsthm}
\usepackage{fancyhdr}
\usepackage{lastpage}
\usepackage{enumitem}
\usepackage{titlesec}
\usepackage{xcolor}
\usepackage{tikz}
\usepackage{hyperref}
\usepackage{bookmark}
\usepackage[most]{tcolorbox}
\usepackage{setspace}
\usepackage{array}
\usetikzlibrary{arrows.meta,positioning}

\hypersetup{
  colorlinks=true,
  linkcolor=blue!50!black,
  urlcolor=blue!50!black,
  citecolor=blue!50!black
}

\setstretch{1.15}
\setlength{\parindent}{2em}
\setlength{\parskip}{0.28em}
\allowdisplaybreaks
\setlist[enumerate]{itemsep=0.35em, topsep=0.35em, leftmargin=2.2em}
\setlist[itemize]{itemsep=0.25em, topsep=0.25em, leftmargin=1.9em}

\definecolor{mainblue}{RGB}{36,83,140}
\definecolor{lightblue}{RGB}{241,246,252}
\definecolor{lightgray}{RGB}{246,246,246}
\definecolor{accent}{RGB}{153,76,0}
\definecolor{rulegray}{RGB}{110,118,128}

\titleformat{\section}
  {\Large\bfseries\color{mainblue}}
  {\thesection.}
  {0.45em}
  {}
  [\vspace{0.25em}\color{rulegray}\titlerule]
\titlespacing*{\section}{0pt}{1.2em}{0.8em}
\titleformat{\subsection}{\large\bfseries\color{mainblue}}{\thesubsection}{0.5em}{}

\pagestyle{fancy}
\fancyhf{}
\fancyhead[L]{\small 欧拉路径与哈密顿路径提高训练题单}
\fancyhead[C]{\small 组合专题：图上走法与结构判断}
\fancyhead[R]{\small 刘文博}
\fancyfoot[C]{\small 第 \thepage 页 / 共 \pageref{LastPage} 页}
\renewcommand{\headrulewidth}{0.55pt}
\renewcommand{\footrulewidth}{0.35pt}

\newtcolorbox{goalbox}[1]{
  breakable,
  colback=lightblue,
  colframe=mainblue,
  boxrule=0.7pt,
  arc=1mm,
  title=\textbf{#1},
  fonttitle=\bfseries
}

\newtcolorbox{notebox}[1]{
  breakable,
  colback=lightgray,
  colframe=black!50,
  boxrule=0.55pt,
  arc=1mm,
  title=\textbf{#1},
  fonttitle=\bfseries
}

\newtheoremstyle{formalexample}
  {0.8em}{0.8em}{\normalfont}{}{\bfseries\color{accent}}{．}{0.6em}{}
\theoremstyle{formalexample}
\newtheorem{exercise}{题目}[section]
\newenvironment{solution}{\par\noindent\textbf{解：}}{\hfill$\square$\par}

\begin{document}

\thispagestyle{empty}
\begin{titlepage}
\centering
\vspace*{0.9cm}
\begin{tcolorbox}[
  enhanced,
  width=0.92\textwidth,
  colback=white,
  colframe=mainblue,
  boxrule=1pt,
  arc=0mm,
  left=6mm,
  right=6mm,
  top=8mm,
  bottom=8mm
]
\centering
{\fontsize{24}{30}\selectfont\bfseries 欧拉路径与哈密顿路径提高训练题单\par}
\vspace{0.6cm}
{\fontsize{17}{22}\selectfont 适合小学高年级至初中低年级竞赛过渡训练\par}

\vspace{1cm}
\begin{tcolorbox}[colback=lightblue,colframe=mainblue,width=0.84\textwidth,boxrule=1pt]
\centering
{\large 训练目标：把“会区分边一次和点一次”提升到“会判断、会构造、会排除”}\\[0.35em]
{\normalsize 建议分 2--3 次完成，先独立做题，再对照详细解答订正}
\end{tcolorbox}

\vspace{1cm}
\renewcommand{\arraystretch}{1.42}
\begin{tabular}{>{\bfseries}p{3.5cm}p{8cm}}
题目数量： & 12 题，按难度递进 \\
专题重点： & 欧拉道路、欧拉回路、奇度点、哈密顿路径、哈密顿回路、二分图染色、结构排除 \\
答案形式： & 末尾附较详细解答，强调“这是边的问题还是点的问题”“为什么必要条件还不够” \\
编译方式： & XeLaTeX \\
\end{tabular}

\vspace{0.9cm}
{\large \today\par}
\end{tcolorbox}
\end{titlepage}

\tableofcontents
\newpage

\section{使用建议}

\begin{goalbox}{做这类题时先问自己的 4 个问题}
\begin{itemize}
  \item 题目要求的是每条边一次，还是每个点一次？
  \item 如果是欧拉题，图是否连通，奇度点有几个？
  \item 如果是哈密顿题，能不能先用度数为 $1$ 的点、二分图染色、分割结构去排除？
  \item 如果必要条件都没问题，能否自己画出一条具体走法？
\end{itemize}
\end{goalbox}

\begin{notebox}{建议顺序}
\begin{itemize}
  \item 第 1--4 题先稳住欧拉判定；
  \item 第 5--8 题训练哈密顿路径与回路的必要条件；
  \item 第 9--12 题体会“必要条件不等于充分条件”和图论结构分析。
\end{itemize}
\end{notebox}

\section{欧拉道路与欧拉回路}

\begin{center}
\begin{tikzpicture}[
  scale=0.85,
  every node/.style={circle,draw=mainblue,fill=lightblue,inner sep=1.5pt,font=\small},
  edge/.style={draw=black!65,line width=0.9pt},
  odd/.style={circle,draw=accent,fill=orange!18,inner sep=1.5pt,font=\small\bfseries}
]
  \node[draw=none,fill=none,circle=false,font=\bfseries\color{mainblue}] at (1.4,3.4) {$K_4$};
  \node (A) at (0,1.2) {};
  \node (B) at (1.4,2.6) {};
  \node (C) at (2.8,1.2) {};
  \node (D) at (1.4,-0.2) {};
  \foreach \u/\v in {A/B,B/C,C/D,D/A,A/C,B/D} {\draw[edge] (\u)--(\v);}
  \node[draw=none,fill=none,circle=false,font=\scriptsize] at (1.4,-1.0) {4 个点度数都为 3};

  \begin{scope}[xshift=5.8cm]
    \node[draw=none,fill=none,circle=false,font=\bfseries\color{mainblue}] at (1.4,3.4) {共点双三角形};
    \node[odd] (O) at (1.4,1.1) {};
    \node (P) at (0.2,2.3) {};
    \node (Q) at (0.2,-0.1) {};
    \node (R) at (2.6,2.3) {};
    \node (S) at (2.6,-0.1) {};
    \draw[edge] (O)--(P)--(Q)--(O);
    \draw[edge] (O)--(R)--(S)--(O);
    \node[draw=none,fill=none,circle=false,font=\scriptsize] at (1.4,-1.0) {公共点度数 4，其余点度数 2};
  \end{scope}

  \begin{scope}[xshift=11.6cm]
    \node[draw=none,fill=none,circle=false,font=\bfseries\color{mainblue}] at (1.5,3.4) {$2\times 3$ 网格图};
    \foreach \x in {0,1,2,3} {
      \foreach \y in {0,1,2} {
        \pgfmathtruncatemacro{\isodd}{ifthenelse((\x==1 || \x==2) && (\y==0 || \y==2),1,0)}
        \ifnum\isodd=1
          \node[odd] (\x\y) at (0.9*\x,0.9*\y) {};
        \else
          \node (\x\y) at (0.9*\x,0.9*\y) {};
        \fi
      }
    }
    \foreach \u/\v in {00/10,10/20,20/30,01/11,11/21,21/31,02/12,12/22,22/32,00/01,01/02,10/11,11/12,20/21,21/22,30/31,31/32} {
      \draw[edge] (\u)--(\v);
    }
    \node[draw=none,fill=none,circle=false,font=\scriptsize] at (1.35,-1.0) {4 个奇度点};
  \end{scope}
\end{tikzpicture}
\end{center}

\begin{exercise}
一个连通图有 $6$ 个奇度点。判断它是否可能存在欧拉道路，并说明理由。
\end{exercise}

\begin{exercise}
正方形 $ABCD$ 的对角线 $AC$ 和 $BD$ 都画出。这个图有没有欧拉回路？有没有欧拉道路？
\end{exercise}

\begin{exercise}
一个图由两个三角形共用一个顶点组成。判断这个图有没有欧拉回路，并说明理由。
\end{exercise}

\begin{exercise}
把 $2\times 3$ 的小方格网格看成图：格点是顶点，相邻格点连边。判断这个图有没有欧拉道路。
\end{exercise}

\section{哈密顿路径与哈密顿回路}

\begin{center}
\begin{tikzpicture}[
  scale=0.85,
  every node/.style={circle,draw=mainblue,fill=lightblue,inner sep=1.5pt,font=\small},
  edge/.style={draw=black!65,line width=0.9pt},
  cycle/.style={draw=green!45!black,line width=1.2pt,-{Stealth[length=2mm]}}
]
  \node[draw=none,fill=none,circle=false,font=\bfseries\color{mainblue}] at (0,3.5) {星形图 $K_{1,5}$};
  \node[fill=orange!18,draw=accent] (O) at (0,1.2) {};
  \foreach \name/\x/\y in {A/-1.8/2.6,B/-2.1/0.8,C/-1.4/-0.8,D/1.6/2.5,E/2.0/0.3} {
    \node (\name) at (\x,\y) {};
    \draw[edge] (O)--(\name);
  }
  \node[draw=none,fill=none,circle=false,font=\scriptsize] at (0,-1.7) {叶子太多，端点装不下};

  \begin{scope}[xshift=6.0cm]
    \node[draw=none,fill=none,circle=false,font=\bfseries\color{mainblue}] at (1.5,3.5) {$3\times 3$ 网格图};
    \foreach \x in {0,1,2} {
      \foreach \y in {0,1,2} {
        \pgfmathtruncatemacro{\c}{mod(\x+\y,2)}
        \ifnum\c=0
          \node[fill=black!70,draw=black!70] (\x\y) at (1.2*\x,1.0*\y) {};
        \else
          \node[fill=white,draw=black!70] (\x\y) at (1.2*\x,1.0*\y) {};
        \fi
      }
    }
    \foreach \x/\y in {00/10,10/20,01/11,11/21,02/12,12/22,00/01,01/02,10/11,11/12,20/21,21/22} {
      \draw[edge] (\x)--(\y);
    }
    \node[draw=none,fill=none,circle=false,font=\scriptsize] at (1.2,-1.0) {黑 5、白 4，无法回路交替};
  \end{scope}

  \begin{scope}[xshift=11.7cm]
    \node[draw=none,fill=none,circle=false,font=\bfseries\color{mainblue}] at (1.8,3.5) {$2\times 4$ 网格图};
    \foreach \x in {0,1,2,3,4} {
      \foreach \y in {0,1} {
        \node (\x\y) at (0.9*\x,1.0*\y) {};
      }
    }
    \foreach \u/\v in {00/10,10/20,20/30,30/40,01/11,11/21,21/31,31/41,00/01,10/11,20/21,30/31,40/41} {
      \draw[edge] (\u)--(\v);
    }
    \draw[cycle] (00)--(10)--(20)--(30)--(40)--(41)--(31)--(21)--(11)--(01)--(00);
    \node[draw=none,fill=none,circle=false,font=\scriptsize] at (1.8,-1.0) {外圈正好经过全部 10 个点};
  \end{scope}
\end{tikzpicture}
\end{center}

\begin{exercise}
星形图 $K_{1,5}$ 有没有哈密顿路径？有没有哈密顿回路？
\end{exercise}

\begin{exercise}
$3\times 3$ 网格图有没有哈密顿回路？请用染色说明。
\end{exercise}

\begin{exercise}
$2\times 4$ 网格图有没有哈密顿回路？如果有，请说明为什么；如果没有，请说明原因。
\end{exercise}

\begin{exercise}
一个图有 7 个顶点，其中有 3 个顶点的度数都是 $1$。证明这个图不可能有哈密顿路径。
\end{exercise}

\section{综合判断与构造}

\begin{exercise}
正五边形的五条边和全部五条对角线组成完全图 $K_5$。判断它有没有欧拉回路、有没有哈密顿回路。
\end{exercise}

\begin{exercise}
请举出一个“有欧拉回路但没有哈密顿回路”的图，并说明理由。
\end{exercise}

\begin{exercise}
请举出一个“有哈密顿回路但没有欧拉道路”的图，并说明理由。
\end{exercise}

\begin{exercise}
为什么说“每个点度数至少为 $2$”只是哈密顿回路的必要条件，不是充分条件？请结合一个具体图形说明。
\end{exercise}

\newpage
\section{详细解答}

\begin{enumerate}
  \item \begin{solution}
  不可能。

  因为欧拉问题的判定非常明确：
  \begin{itemize}
    \item 若有欧拉回路，则奇度点个数必须是 $0$；
    \item 若有欧拉道路但没有欧拉回路，则奇度点个数必须是 $2$。
  \end{itemize}

  题目给出的连通图有 $6$ 个奇度点，既不是 $0$，也不是 $2$，所以它既不可能有欧拉回路，也不可能有欧拉道路。
  \end{solution}

  \item \begin{solution}
  先数度数。

  在这个图里，每个顶点都和另外三个顶点相连，所以四个顶点的度数都是 $3$。

  因此图中有 $4$ 个奇度点。

  欧拉回路要求所有点度数都为偶数，所以这里没有欧拉回路。

  欧拉道路要求恰好有 $2$ 个奇度点，而现在有 $4$ 个奇度点，所以这里也没有欧拉道路。
  \end{solution}

  \item \begin{solution}
  有欧拉回路。

  设两个三角形共用的顶点为 $O$。那么：
  \begin{itemize}
    \item 每个三角形里除了 $O$ 之外的两个顶点，度数都是 $2$；
    \item 公共顶点 $O$ 在两个三角形中各连出两条边，所以度数是 $4$。
  \end{itemize}

  这样一来，所有顶点度数都是偶数，而且整个图是连通的。

  由欧拉回路判定定理，这个图有欧拉回路。

  例如，可以先绕第一个三角形一圈回到 $O$，再绕第二个三角形一圈回到 $O$。
  \end{solution}

  \item \begin{solution}
  没有欧拉道路。

  $2\times 3$ 小方格网格有 $3\times 4=12$ 个格点。观察边界上的点和中间点可知：
  \begin{itemize}
    \item 四个角点的度数是 $2$；
    \item 边上但不是角的点，度数是 $3$；
    \item 内部点的度数是 $4$。
  \end{itemize}

  在这个图里，上边中间两个点和下边中间两个点是度数 $3$ 的点，共有 $4$ 个奇度点。

  由于奇度点个数是 $4$，不等于 $0$ 或 $2$，所以不存在欧拉道路。
  \end{solution}

  \item \begin{solution}
  都没有。

  星形图 $K_{1,5}$ 中有一个中心点和 $5$ 个叶子点。叶子点之间没有边。

  先看哈密顿回路：回路要求每个点都要在回路上“进一次、出一次”，因此回路中的每个点度数至少应该为 $2$。但这里每个叶子点度数都是 $1$，所以根本不可能有哈密顿回路。

  再看哈密顿路径：一条路径最多只有两个端点，所以最多只能容纳两个度数为 $1$ 的叶子点作为路径端点。现在却有 $5$ 个叶子点，显然不可能全部放进同一条哈密顿路径中。

  因此它既没有哈密顿路径，也没有哈密顿回路。
  \end{solution}

  \item \begin{solution}
  没有哈密顿回路。

  把 $3\times 3$ 网格图做黑白染色。因为共有 $9$ 个点，所以两色点数分别是 $5$ 个和 $4$ 个。

  网格图是二分图，任何回路上的相邻点颜色都必须交替。因此若存在哈密顿回路，那么回路经过的黑点数和白点数必须相等。

  但现在黑白点数不相等，所以不可能存在哈密顿回路。
  \end{solution}

  \item \begin{solution}
  有哈密顿回路。

  把 $2\times 4$ 网格图看成两行五列的格点图。最外层这些点围成一个大长方形边界。

  沿着这个边界走一圈，就可以恰好经过所有边界上的点一次并回到起点；而在 $2\times 4$ 网格图中，所有点都恰好都在这个边界上，所以这条边界回路正是一条哈密顿回路。

  因此它有哈密顿回路。
  \end{solution}

  \item \begin{solution}
  设这 3 个度数为 $1$ 的顶点分别是 $A,B,C$。

  如果图有哈密顿路径，那么这条路径要经过每个顶点一次。路径中只有两个端点，端点可以只连出一条被使用的边；而路径中间的点必须至少连着两条被使用的边，一条进入，一条离开。

  度数为 $1$ 的顶点不可能放在路径中间，因此 $A,B,C$ 这三个点都只能当路径端点。

  但是一条路径最多只有两个端点，不可能同时容纳三个度数为 $1$ 的顶点。

  矛盾，所以该图不可能有哈密顿路径。
  \end{solution}

  \item \begin{solution}
  完全图 $K_5$ 中每个顶点都与其余 $4$ 个顶点相连，所以每个顶点度数都是 $4$。

  先看欧拉：图连通，且所有顶点度数都是偶数，因此有欧拉回路。

  再看哈密顿：完全图里任意两个不同顶点之间都有边，当然可以依次经过五个顶点再回到起点，所以有哈密顿回路。

  因此 $K_5$ 既有欧拉回路，也有哈密顿回路。
  \end{solution}

  \item \begin{solution}
  一个典型例子，就是“两只三角形只共用一个顶点”的图。

  上面第 3 题已经说明，这个图所有顶点度数都是偶数，因此它有欧拉回路。

  但它没有哈密顿回路。原因是：如果想从一个三角形走到另一个三角形，必须经过公共顶点；而如果想绕完两个三角形再回到出发点，公共顶点就必须被经过不止一次。

  可哈密顿回路要求每个顶点只能经过一次，所以这是不可能的。

  于是这个图就是“有欧拉回路但没有哈密顿回路”的例子。
  \end{solution}

  \item \begin{solution}
  一个简单例子是完全图 $K_4$。

  在 $K_4$ 中，每个顶点都与其余 $3$ 个顶点相连，所以每个顶点度数都是 $3$。

  因为有 $4$ 个奇度点，所以它没有欧拉道路。

  但是它有哈密顿回路。例如按四个顶点依次绕一圈，就能每个顶点经过一次再回到起点。

  所以 $K_4$ 是“有哈密顿回路但没有欧拉道路”的一个例子。
  \end{solution}

  \item \begin{solution}
  因为“每个点度数至少为 $2$”只说明每个点局部上看起来有希望进出一次，但并没有保证整个图能用一条回路把所有点串起来。

  一个具体例子仍然可以用“两只三角形共用一个顶点”的图。

  在这个图中：
  \begin{itemize}
    \item 公共顶点度数是 $4$；
    \item 其余四个顶点度数都是 $2$。
  \end{itemize}

  所以每个点度数都至少为 $2$，必要条件完全满足。

  但它还是没有哈密顿回路，因为如果想把两个三角形都绕到，就必须两次经过公共顶点，这违反了哈密顿回路“每个点只经过一次”的要求。

  因此，“每个点度数至少为 $2$”只是必要条件，不是充分条件。
  \end{solution}
\end{enumerate}

\end{document}
